%! Author = matteomagnini
%! Date = 22/04/26

\documentclass[11pt]{article}

\usepackage{amsmath}
\usepackage{amssymb}
\usepackage{my-acronyms}
\usepackage{hyperref}

\title{
    Intelligent System 2\\
    Tutorial week 9: Default Logics
}

\author{
    Matteo Magnini\\
    \href{mailto:matteo.magnini@uni.lu}{matteo.magnini@uni.lu}
}

\date{
    23 April, 2026
}

\begin{document}

\maketitle

\section{Exercises}

\subsection{Exercise 1: Encoding a Default Theory}

Consider the following scenario:

\begin{itemize}
    \item Typically, students submit assignments on time.
    \item Typically, busy people do not submit assignments on time.
    \item Alice is a student.
    \item Alice is busy.
\end{itemize}

Let the propositional variables be:
\[
student, \quad busy, \quad onTime
\]

\begin{enumerate}
    \item Define the default theory $T = (F, D)$:
    \begin{itemize}
        \item Specify the set of facts $F$.
        \item Specify the set of defaults $D$.
    \end{itemize}

    \item Assign a strength function $\text{str}: D \rightarrow \mathbb{N}$ assuming:
    \[
    \text{general traits} < \text{behavioral conditions}
    \]

    \item Intuitively: should Alice submit on time? Explain informally.
\end{enumerate}

\subsection{Exercise 2: Greedy Method (Step-by-Step)}

Consider the default theory:

\[
F = \{ tired \}
\]

\[
D =
\begin{cases}
d_1: \; \top \Rightarrow worker \\
d_2: \; worker \Rightarrow productive \\
d_3: \; tired \Rightarrow \neg productive
\end{cases}
\]

with strengths:
\[
\text{str}(d_1) = 1, \quad \text{str}(d_2) = 2, \quad \text{str}(d_3) = 3
\]

\begin{enumerate}
    \item Compute the Greedy construction:
    \begin{itemize}
        \item Step 0: $gr_0 = \emptyset$
        \item At each step:
        \begin{itemize}
            \item list applicable defaults
            \item select the strongest consistent one
        \end{itemize}
    \end{itemize}

    \item Fill the table:

    \begin{center}
    \begin{tabular}{c|c|c}
        step & $Cn(F \cup hd[gr_n])$ & selected default \\
        \hline
        0 & & \\
        1 & & \\
        2 & & \\
    \end{tabular}
    \end{center}

    \item Give the final Greedy extension.
\end{enumerate}

\subsection{Exercise 3: Hansen Method}

Use the same theory as in Exercise 2.

\begin{enumerate}
    \item Compute the Hansen construction:
    \begin{itemize}
        \item At each step, select the strongest \textbf{consistent} default
        \item Note: it might not be applicable
    \end{itemize}

    \item Fill the table:

    \begin{center}
    \begin{tabular}{c|c|c}
        step & $out(F, D_n)$ & selected default \\
        \hline
        0 & & \\
        1 & & \\
        2 & & \\
    \end{tabular}
    \end{center}

    \item Which defaults are selected but never activated?

    \item Compare the Hansen extension with the Greedy one.
\end{enumerate}

\subsection{Exercise 4: Brewka-Eiter Method}

Consider again the theory from Exercise 2.

\begin{enumerate}
    \item Take the Greedy extension $S$.

    \item Verify that $S$ is a Greedy extension of $T$.

    \item Compute the reduction $T^S = (F \cup S, D^S, str^S)$:
    \begin{itemize}
        \item Define $D^S$
        \item Assign strengths $str^S$
    \end{itemize}

    \item Compute the Greedy construction for $T^S$.

    \item Is $S$ a Brewka-Eiter extension? Justify.
\end{enumerate}

\subsection{Exercise 5: Full Problem}

Consider the following scenario:

\begin{itemize}
    \item Typically, artists are creative.
    \item Typically, engineers are not creative.
    \item Typically, creative people enjoy abstract art.
    \item Bob is an artist.
    \item Bob is an engineer.
\end{itemize}

Let:
\[
artist, \quad engineer, \quad creative, \quad abstract
\]

\begin{enumerate}
    \item Define the default theory $T = (F, D, str)$ assuming:
    \[
    \text{profession stereotypes} < \text{cognitive traits}
    \]

    \item Compute the Greedy extension:
    \begin{itemize}
        \item show all steps
    \end{itemize}

    \item Compute the Hansen extension.

    \item Guess a Brewka-Eiter candidate $S$.

    \item Verify:
    \begin{enumerate}
        \item $S$ is Greedy for $T$
        \item $S$ is Greedy for $T^S$
    \end{enumerate}

    \item Final question:
    \begin{quote}
        Does Bob enjoy abstract art?
    \end{quote}
    Justify using the different methods.
\end{enumerate}

\subsection{Exercise 6 (Optional, Conceptual)}

Answer briefly:

\begin{enumerate}
    \item Why does the Greedy method sometimes produce more conclusions than Hansen?

    \item Why can Hansen discard useful defaults?

    \item What is the intuition behind Brewka-Eiter fixing this issue?

\end{enumerate}

\end{document}